\chapter{LightGBM and XGBoost}
\label{ch:lgbm-xgb}

\begin{saybox}[title=What to say at the start of this section]
Explain boosting as correction rather than averaging: each tree is fitted
to what the previous trees got wrong. Then make the key comparison
explicit --- under the original protocol the three ensembles were
indistinguishable; under leave-rule-out they separate, and they fail on
the same rules.
\end{saybox}

\section{Gradient boosting}

Both models fit an additive ensemble stage by stage,
\begin{equation}
F_M(\mathbf{x}) = F_0(\mathbf{x}) + \eta \sum_{m=1}^{M} h_m(\mathbf{x}),
\end{equation}
where each $h_m$ is a regression tree fitted to the negative gradient of
the loss at the current prediction. For squared error the gradient is the
residual itself,
\begin{equation}
r_i^{(m)} = -\left. \frac{\partial \ell(y_i, F)}{\partial F} \right|_{F = F_{m-1}(\mathbf{x}_i)}
= y_i - F_{m-1}(\mathbf{x}_i),
\end{equation}
so each stage literally models what the ensemble has so far failed to
explain. The learning rate $\eta$ shrinks each stage's contribution;
small $\eta$ with large $M$ is the standard bias--variance compromise, and
both production models use $\eta = 0.03$ with $M \le 2000$ under early
stopping.

\subsection{XGBoost's regularised objective}

XGBoost optimises a second-order approximation with an explicit complexity
penalty on each tree:
\begin{equation}
\mathcal{L}^{(m)} \simeq \sum_{i=1}^{n}
\left[ g_i h_m(\mathbf{x}_i) + \tfrac{1}{2} \tilde{h}_i h_m(\mathbf{x}_i)^2 \right]
+ \gamma T + \tfrac{1}{2}\lambda \sum_{j=1}^{T} w_j^2 ,
\end{equation}
with $T$ the number of leaves and $w_j$ the leaf weights. The optimal leaf
weight is $w_j^\star = -G_j / (H_j + \lambda)$, so the $L_2$ term shrinks
leaves supported by few rows --- which is precisely the situation in the
small composite rules of \Cref{sec:fold-natamycin}.

\subsection{Leaf-wise versus depth-wise growth}

LightGBM grows leaf-wise: at each step it splits the leaf offering the
largest loss reduction anywhere in the tree, producing deep, unbalanced
trees that chase the hardest regions. XGBoost grows depth-wise to a fixed
\texttt{max\_depth}, producing balanced trees. LightGBM's strategy is
typically more sample-efficient and more prone to over-specialising on
small regions; \texttt{num\_leaves=31} and
\texttt{min\_child\_samples=20} are the two settings that bound it here.

\section{Production configurations}

\begin{center}\small
\begin{tabular}{lp{5.3cm}p{5.3cm}}
\toprule
& LightGBM & XGBoost \\
\midrule
Rounds & \texttt{n\_estimators=2000}, early stopping 50 &
\texttt{n\_estimators=2000}, \texttt{early\_stopping\_rounds=50} \\
Learning rate & \texttt{0.03} & \texttt{0.03} \\
Tree shape & \texttt{num\_leaves=31} (leaf-wise) & \texttt{max\_depth=6} (depth-wise) \\
Leaf constraint & \texttt{min\_child\_samples=20} & default $\lambda$, $\gamma$ \\
Sampling & \texttt{subsample=0.9}, \texttt{colsample\_bytree=0.9} & identical \\
\bottomrule
\end{tabular}
\end{center}

Early stopping monitors the validation partition, which means that under
the original protocol the stopping iteration was selected using data drawn
from the same generation rules as the test set. In the leave-rule-out
harness, boosting rounds are fixed a priori rather than early-stopped,
which removes that dependency at the cost of not tuning $M$ per fold. The
difference is documented rather than silently absorbed.

\section{Original results}

From \Cref{tab:tab_production_v7}: LightGBM is the best model on five of
the six V7 specialists under the original protocol, reaching $0.971$
(soft/general), $0.920$ (semi-hard/general) and $0.932$ (hard/general).
XGBoost is within one or two points everywhere. The gap between them, and
between either and Random Forest, is well inside the range that a different
random seed could produce --- which is itself a hint that the protocol was
not discriminating between models.

\section{Revised results}

\begin{table}[htbp]\centering\small
\caption{LightGBM and XGBoost under leave-one-rule-out, per specialist.
Computed from the $168$ reconstructed fold results.}
\label{tab:lgbm_xgb_loro}
\begin{tabular}{llrrrrr}
\toprule
Specialist & Model & folds & mean $\rsq$ & median $\rsq$ & min $\rsq$ & pooled OOF $\rsq$ \\
\midrule
\multirow{2}{*}{soft / general} & LightGBM & 7 & 0.850 & 0.892 & 0.472 & 0.960 \\
 & XGBoost & 7 & 0.651 & 0.806 & $-0.279$ & 0.957 \\
\addlinespace[2pt]
\multirow{2}{*}{semi-hard / general} & LightGBM & 12 & 0.824 & 0.888 & 0.562 & 0.940 \\
 & XGBoost & 12 & 0.748 & 0.840 & 0.270 & 0.918 \\
\addlinespace[2pt]
\multirow{2}{*}{hard / general} & LightGBM & 8 & 0.256 & 0.714 & $-2.665$ & 0.669 \\
 & XGBoost & 8 & $-0.033$ & 0.568 & $-2.305$ & 0.665 \\
\addlinespace[2pt]
\multirow{2}{*}{soft / safety} & LightGBM & 11 & 0.702 & 0.831 & $-0.256$ & 0.867 \\
 & XGBoost & 11 & 0.589 & 0.671 & $-0.164$ & 0.786 \\
\addlinespace[2pt]
\multirow{2}{*}{semi-hard / safety} & LightGBM & 8 & 0.860 & 0.926 & 0.617 & 0.811 \\
 & XGBoost & 8 & 0.649 & 0.690 & 0.269 & 0.655 \\
\addlinespace[2pt]
\multirow{2}{*}{hard / safety} & LightGBM & 10 & 0.839 & 0.934 & 0.425 & 0.730 \\
 & XGBoost & 10 & 0.697 & 0.736 & 0.366 & 0.593 \\
\bottomrule
\end{tabular}
\end{table}

\tbllong{tab_folds_lightgbm}{Complete LightGBM fold results for every held-out
generation rule.}

\tbllong{tab_folds_xgboost}{Complete XGBoost fold results for every held-out
generation rule.}

\begin{findingbox}[title=LightGBM is the strongest model under the revised protocol]
LightGBM has the highest median fold $\rsq$ in five of six specialists and
the highest pooled out-of-fold $\rsq$ in all six. Its advantage is largest
exactly where the problem is hardest: on hard/general it reaches mean
$\rsq = 0.256$ against XGBoost's $-0.033$ and Random Forest's $-0.097$,
and on the clostridia fold it is the least badly wrong of the three
($-2.665$ against $-2.305$ and $-3.860$; MAE $55.1$ days against $57.3$
and $67.4$).

This ordering was not visible under the original protocol, where all three
sat within two points of each other. An evaluation that cannot distinguish
between candidate models cannot be used to select between them.
\end{findingbox}

\section{Do the three ensembles fail on the same rules?}

\fig{fig34_model_agreement}{0.98}{%
  Pairwise per-fold $\rsq$ for the three ensembles across all $56$ folds,
  clipped at $-3$. Correlations: Random Forest--LightGBM $r = 0.939$,
  Random Forest--XGBoost $r = 0.877$, LightGBM--XGBoost $r = 0.878$.}

The agreement is high. All three models place the clostridia fold as their
single worst result, and all three struggle on the shredded-cheddar rule.
The failures are shared, not model-specific.

\begin{findingbox}[title=What correlated failure does and does not imply]
Correlation of $r \approx 0.9$ across three structurally different
learners --- one bagging, two boosting, one depth-wise and one leaf-wise
--- indicates that fold difficulty is a property of the \emph{data
partition}, not of any model's inductive bias.

The careful non-claim: this does not by itself prove the cause is
mechanism novelty. Three tree ensembles share the extrapolation bound of
\cref{eq:rf-bound}, so they would be expected to fail together on any
out-of-range target regardless of why the target is out of range. The
causal attribution in \Cref{sec:fold-clostridia} rests on the two
controlled comparisons, not on this correlation.
\end{findingbox}

\fig{fig35_rule_model_heatmap_r2}{0.62}{%
  Per-rule $\rsq$ for all three models across all $56$ held-out rules,
  sorted by mean difficulty.}

\fig{fig36_rule_model_heatmap_mae}{0.62}{%
  The same grid in MAE (days), which removes the variance-denominator
  effect. Several rules that look poor in $\rsq$ have small day-scale
  errors.}
